\section{Introduction}\label{sec:intro}
Can every class learned by statistical queries under arbitrary input distributions also be learned in a low-dimensional feature space? A positive answer would make fixed linear representations universal for this form of learning. The requirement is not automatic: a learner may use its answers to decide where the distribution places mass, whereas a common feature map must be chosen before the distribution is known. This is Open Question~2 of \citet{FKS26}.

We give a negative answer with an exponential separation. For infinitely many input lengths $n$, our classes admit $O_\eps(n)$ queries of constant tolerance, but their ordinary, exact probabilistic, and expected-error dimensions are all $\Theta(2^{n/3})$ at fixed approximation error below $1/2$. The lower bound survives a further relaxation: the feature law may depend on a known target prior and may fail on any fixed fraction of targets below one. The learners can be proper and deterministic. Their computation is polynomial in the explicit table, not in $n$.

\citet{Patel26} posted the first public separation. His read-once DNFs have ordinary dimension $2^{\Omega(n^{1/3})}$ and, at error $0.01$, an SQ learner using $2^{\widetilde O(n^{2/9})}$ queries of tolerance $2^{-\widetilde O(n^{2/9})}$. Our bounds improve the query, tolerance, and dimension parameters and cover probabilistic representations, but the construction is a random table rather than a natural succinct class. The qualitative separation also follows from \citet{DKR25} and \citet{RS10}, as Patel makes explicit.
\ifreview
This work was obtained independently of \citet{Patel26}.
\else
This work was obtained independently of \citet{Patel26}; a draft containing the main separation was shared with V.\ Feldman on 11 September 2026, before Patel's preprint appeared.
\fi

\subsection{Proof overview}\label{sec:techniques}
We use the incidence construction of \citet{HHPTZ22}. Rows index integer lines and columns index grid points; the template is negative exactly on incidences. Give these incidences arbitrary signs, keeping all other entries positive. Every resulting table still has large monochromatic rectangles under every product distribution, and two rows still share at most one negative point. Choosing the incidence signs independently supplies enough randomness to force high dimension.

The learner uses a rectangle in one of two ways. If the target rarely disagrees with its color, label the rectangle's points and remove them. Otherwise, remove its hypotheses. The first operation reduces the mass left to label; the second reduces the set of possible targets. A minimax argument supplies a distribution over rectangles without knowing the point marginal. Either accepted update makes progress at least equal to the product of the rectangle's two masses. This common lower bound controls the query count despite adversarial oracle answers.

The dimension argument counts only the entries whose signs vary. At scale $N$, these are $\Theta(N^4)$ incidences in a table with $\Theta(N^3)$ rows and columns. Warren's theorem gives $2^{O(N^3d\log(N/d))}$ sign patterns on this support when $d$ is a sufficiently small fraction of $N$. Taking $d=cN$ with small constant $c$ leaves too few patterns to cover the random signs, even after allowing a fixed fraction of errors. This removes the logarithmic loss in the dimension bound of \citet{HHPTZ22}. A union bound over all large row subsets then permits the representation to discard targets.

\subsection{Main result}\label{sec:main}
Write $\dc$ for ordinary dimension, $\edc$ for exact probabilistic dimension with success probability $1/2$, and $\dc_\eta$ for expected-error dimension at error $\eta$. Section~\ref{sec:definitions} gives the definitions. All learners below must succeed against every valid SQ oracle, without access to the unknown marginal.

\begin{theorem}[Sharp incidence separation]\label{thm:main}
For each integer $N\ge2$, there is a family of classes $\Hc_A$ on $2N^3$ points such that every member has VC dimension at most two and $\dc(\Hc_A)\le N+3$. For every $0<\eps<1/4$, every member has a proper deterministic distribution-independent SQ learner with
\begin{equation}
 m=O\bigl(\log N+\eps^{-1}\bigr),
 \qquad \tau=\Omega\bigl(\eps/\log(1/\eps)\bigr).
 \label{eq:mainlearn}
\end{equation}
The deterministic improper learner uses only $O(\log N+\log(1/\eps))$ queries at the same tolerance; the $1/\eps$ term comes from the row tests in proper completion. Both learners run in time polynomial in the explicit table size and $1/\eps$. For each fixed $0\le\eta<1/2$, independent fair incidence signs give, with probability $1-\exp(-\Omega_\eta(N^4))$,
\begin{equation}
 \dc(\Hc_A)=\Theta(N),\qquad
 \edc(\Hc_A)=\Theta(N),\qquad
 \dc_\eta(\Hc_A)=\Theta_\eta(N).
 \label{eq:maindim}
\end{equation}
For each fixed $0\le\eta<1/2$ and $0\le\delta<1$, the uniform prior on the full-length lines also has prior-average dimension $\Theta_{\eta,\delta}(N)$ with probability $1-\exp(-\Omega_{\eta,\delta}(N^4))$.
\end{theorem}

\begin{proofsketch}
The incidence template has sign-rank at most six, which gives a constant rectangle ratio after every flip. Theorem~\ref{thm:rectangle} supplies the learner; Appendix~\ref{app:proper} makes it deterministic and proper.

For dimension, test each row on its incidence support. Any common random feature law of small expected error yields one embedding with few total incidence mistakes. Lemma~\ref{lem:mask} bounds the number of such low-rank patterns. One constant feature reduces exact probabilistic dimension to this case. Conversely, $N$ indicator features and three quadratic features represent every flip. Applying the count to every sufficiently large row subset gives the prior-average assertion of Theorem~\ref{thm:prior}.
\end{proofsketch}

Taking $N=2^b$ gives a domain of size $2^n$ with $n=3b+1$. Thus no universal polynomial in $(m,1/\tau,n)$ bounds any of the three dimensions. This also refutes a bound on $\dc_{C\eps}$ for any universal $C>0$: choose $\eps$ once so that $C\eps<1/2$, then let $N$ grow. At error $1/2$, a constant feature would suffice.

The query order is also optimal: typical tables require $\Theta_\eps(n)$ queries at a suitable fixed tolerance, even for randomized improper learning. In the other direction, a learner's quantized response tree gives exponential upper bounds on representation dimension. These bounds have the same exponential order in query count on the incidence family.

\section{Learning and representation models}\label{sec:definitions}
Let $X$ be a finite nonempty domain and $\Hc=\{h_1,\ldots,h_K\}\subseteq\{-1,+1\}^{X}$ an indexed class. Repeated rows are allowed. Write $Q=|X|$, $n=\ceil{\log_2 Q}$, and $\Delta(Z)$ for the probability simplex on a finite set $Z$. For a marginal $D\in\Delta(X)$, target $h$, and Boolean predictor $g$, put
\[
 L_{D,h}(g)=\Pr_{z\sim D}[g(z)\ne h(z)].
\]
We write $[k]=\{1,\ldots,k\}$, use $\ln$ for natural logarithms and $\log_2$ for binary logarithms, and leave the base unspecified only in asymptotic bounds.

\paragraph{Statistical queries.}
A query $q:X\times\{-1,+1\}\to[-1,1]$ receives a number within $\tau$ of $\E_{z\sim D}q(z,h(z))$ \citep{Kearns98}. A valid oracle policy may depend on previous queries and answers and on the current query, but not on unrevealed future learner randomness. A depth-$m$ learner uses at most $m$ queries and satisfies
\begin{equation}
 \forall D\ \forall h\in\Hc\ \forall\text{ valid policies }O,
 \qquad \E_\omega L_{D,h}(A^O_\omega)\le\eps.
 \label{eq:learnmodel}
\end{equation}
Here $\omega$ is the learner's independent seed. The algorithm may depend on the class, not on the unknown instance $(D,h)$. A proper learner returns a member of $\Hc$. Computation is uncharged unless stated; finite-bit implementations include answer rounding in the tolerance.

\paragraph{Dimension.}
Ordinary dimension complexity $\dc(\Hc)$ is the least $d$ admitting a map $\phi:X\to\R^d$ and row vectors $w_h$ such that
\begin{equation}
 h(z)\ip{w_h}{\phi(z)}>0\qquad(h\in\Hc,\ z\in X).
 \label{eq:dcdef}
\end{equation}
Factoring the matrix of scores shows that this is the sign-rank of the truth table \citep{BES02}. For probabilistic representations we fix $\sgn(0)=+1$ and write
\[
 \mathcal E_\phi(D,h)=\min_w L_{D,h}(\sgn\ip{w}{\phi}).
\]
The minimum is attained because only finitely many Boolean prediction patterns occur. The optimizing vector may depend on the realized map, the target, and the marginal.

The exact probabilistic dimension $\edc(\Hc)$ is the least $d$ admitting one law on maps $\phi:X\to\R^d$ such that
\begin{equation}
 \forall D,h,\qquad \Pr_\phi[\mathcal E_\phi(D,h)=0]\ge\tfrac12.
 \label{eq:exactprob}
\end{equation}
The expected-error probabilistic dimension $\dc_\eta(\Hc)$ is defined by replacing this condition with
\begin{equation}
 \forall D,h,\qquad \E_\phi\mathcal E_\phi(D,h)\le\eta.
 \label{eq:avgprob}
\end{equation}
In both definitions the feature law is chosen before the target and marginal. We use the probabilistic representation notions of \citet{CMW25} and \citet{KMS20}, with the stated Boolean threshold convention.

\begin{remark}[Threshold ties]\label{rem:ties}
One constant feature turns any finite Boolean prediction matrix into a strict representation: add a row-specific positive bias smaller than every nonzero negative score in absolute value. The opposite fixed tie rule uses a negative bias. Counting zero scores as errors only makes approximation harder. Our lower bounds therefore do not depend on the choice between these conventions, apart from an additive dimension.
\end{remark}

\section{A learner from monochromatic rectangles}\label{sec:rectangle}
The learner needs a region on which many surviving hypotheses agree. For a rectangle $S\times T\subseteq X\times\Hc$, call $c\in\{-1,+1\}$ its color when $h(z)=c$ throughout the rectangle. Following \citet{HHPTZ22}, define
\begin{equation}
 \rect(\Hc)=\inf_{\mu\in\Delta(X),\,\nu\in\Delta(\Hc)}
 \max_{S\times T\text{ monochromatic}}\mu(S)\nu(T).
 \label{eq:rectdef}
\end{equation}
Restrictions to subsets of rows or points cannot decrease this quantity: extend the restricted weights by zero.

\begin{theorem}[Rectangle learner]\label{thm:rectangle}
If $\rect(\Hc)\ge1/\rho$ with $\rho\ge1$, then for every $0<\eps<1/4$ a randomized learner satisfies \eqref{eq:learnmodel} using
\begin{equation}
 m=O\bigl(\rho(\log K+\log(1/\eps))\bigr),
 \qquad \tau=\Omega\bigl(\eps/(\rho\log(1/\eps))\bigr).
 \label{eq:rectanglebound}
\end{equation}
The query budget is a worst-case bound. The learner can be constructed from the table and the stated lower bound on the rectangle ratio.
\end{theorem}

The theorem assumes only a rectangle ratio; finding the best rectangle for the unknown marginal is not part of the learner's input. The following minimax lemma supplies a distribution over candidates instead. This is the minimax argument used by \citet{HHPTZ22} for average margin.

\begin{lemma}[Pointwise coverage]\label{lem:mixture}
Fix nonempty $U\subseteq X$, $V\subseteq\Hc$, a prior $\nu$ on $V$, and $0<r\le\rect(\Hc)$. There is a law $\lambda$ on monochromatic rectangles $S\times T\subseteq U\times V$ satisfying
\begin{equation}
 \E_{(S,T,c)\sim\lambda}[\one_S(z)\nu(T)]\ge r
 \qquad(z\in U).
 \label{eq:pointwise}
\end{equation}
This law depends on the table, $U$, $V$, and $\nu$, but not on $D$.
\end{lemma}

\begin{proof}
For every distribution $\mu$ on $U$, some rectangle has $\mu(S)\nu(T)\ge r$. Finite minimax gives a mixture of these rectangles whose weighted coverage is at least $r$ at every point. Appendix~\ref{app:rectangle} gives the corresponding finite separation argument.
\end{proof}

\paragraph{The algorithm.}\label{alg:rectify}
\Rectify{} maintains unlabelled points $U$, surviving hypotheses $V$, and a partial predictor. Start with $U=X$ and $V=\Hc$. Query $u=D(U)$ and stop when its estimate is at most $\eps/4$. Otherwise use Lemma~\ref{lem:mixture}, with the uniform prior on $V$, to propose $(S,T,c)$. Query $s=D(S)$ and reject the proposal if $\widehat s<(r/2)\widehat u$. For an accepted proposal, query the disagreement mass
\[
 w=D(S\cap\{h\ne c\}).
\]
If $\widehat w\le2\tau$, label $S$ with $c$ and remove it from $U$. Otherwise remove $T$ from $V$. After a fixed number of rounds, stop and extend the partial predictor by $+1$.

\begin{figure}[t]
\centering
\begin{tikzpicture}[x=1cm,y=1cm,>=Stealth,font=\small]
% Each panel has points on the vertical axis and hypotheses on the horizontal axis.
\begin{scope}
  \draw (0,0) rectangle (3.2,1.9);
  \fill[black!16] (.85,.64) rectangle (2.15,1.28);
  \draw (.85,.64) rectangle (2.15,1.28);
  \node at (1.5,.96) {$S\times T$};
  \node at (1.6,2.17) {propose a rectangle};
  \node[rotate=90] at (-.3,.95) {$U$};
  \node at (1.6,-.25) {$V$};
\end{scope}
\draw[->] (3.4,.96) -- (4.05,.96);
\begin{scope}[xshift=4.3cm]
  \draw (0,0) rectangle (3.2,1.9);
  \fill[pattern=north east lines,pattern color=black!45] (0,.64) rectangle (3.2,1.28);
  \draw[dashed] (0,.64)--(3.2,.64) (0,1.28)--(3.2,1.28);
  \node[fill=white,inner sep=2pt] at (1.6,.96) {label $S$};
  \node at (1.6,2.17) {small disagreement};
  \node[align=center] at (1.6,-.5)
     {$D(U')=(1-a)D(U)$\\progress $a$};
\end{scope}
\node at (7.95,.96) {or};
\begin{scope}[xshift=8.4cm]
  \draw (0,0) rectangle (3.2,1.9);
  \fill[pattern=north east lines,pattern color=black!45] (.85,0) rectangle (2.15,1.9);
  \draw[dashed] (.85,0)--(.85,1.9) (2.15,0)--(2.15,1.9);
  \node[fill=white,inner sep=2pt,align=center] at (1.5,.96) {discard\\$T$};
  \node at (1.6,2.17) {witnessed disagreement};
  \node[align=center] at (1.6,-.5)
     {$|V'|=(1-b)|V|$\\progress $b$};
\end{scope}
\node at (5.8,-1.15) {$a= D(S)/D(U),\quad b=\nu(T),\quad\min\{a,b\}\ge ab$};
\end{tikzpicture}
\caption{An accepted proposal either removes a fraction $a$ of the remaining point mass or a fraction $b$ of the surviving hypotheses. The shaded band marks what is removed. Both updates have progress at least $ab$, the quantity controlled by the pointwise rectangle mixture.}
\label{fig:progress}
\end{figure}


\paragraph{Why it stops.}
Put $a=s/u$ and $b=\nu(T)$. With the tolerance chosen in Appendix~\ref{app:rectangle}, acceptance implies $a\ge r/4$, while every proposal with $a\ge3r/4$ is accepted. A labelling step makes at most $3\tau$ error. An elimination never removes the target.

Assign progress $G=a$ to labelling, $G=b$ to elimination, and $G=0$ to rejection. Either accepted update satisfies $G\ge ab$. Pointwise coverage gives $\E[ab\mid\text{history}]\ge r$, so rejection loses at most $3r/4$ and
\begin{equation}
 \E[G\mid\text{history}]\ge r/4.
 \label{eq:drift}
\end{equation}
The cumulative elimination progress is at most $\ln K$; the cumulative labelling progress is at most $\ln(8/\eps)+1$. An exponential potential therefore bounds the probability of exhausting the round budget. The number of labelling steps is $O(\rho\log(1/\eps))$, explaining the tolerance in \eqref{eq:rectanglebound}. The proof in Appendix~\ref{app:rectangle} treats every allowable oracle answer.

\section{Incidence signs and representation dimension}\label{sec:incidence}
At scale $N$, the point set $X$ and line index set $\Lc$ are both $[N]\times[2N^2]$. A point is $p=(x,y)$ and a line is $\ell=(a,b)$, representing $y=ax+b$. Let
\begin{equation}
 F_{\ell,p}=-1\ \Longleftrightarrow\ y=ax+b.
 \label{eq:grid}
\end{equation}
An incidence flip $A$ changes any subset of these negative entries to $+1$ and leaves every other entry positive. Its row class is $\Hc_A$, and $K=Q=2N^3$. A random flip chooses each incidence sign independently and fairly.

The line $\ell=(1,1)$ is a running example. Its template points are $(1,2),\ldots,(N,N+1)$. Their labels in its row can be chosen arbitrarily; the row remains positive elsewhere. A different line shares at most one of these points, and the two rows' signs at that point are separate random bits.

\begin{figure}[t]
\centering
\begin{tikzpicture}[x=1cm,y=1cm,font=\small,>=Stealth]
\begin{scope}[xshift=0cm]
\node at (2.75,2.65) {incidence template $F$};
\node[font=\scriptsize] at (1.1,2.2) {$(1,2)$};
\node[font=\scriptsize] at (1.8900000000000001,2.2) {$(2,3)$};
\node[font=\scriptsize] at (2.68,2.2) {$(3,4)$};
\node[font=\scriptsize] at (3.47,2.2) {$(1,3)$};
\node[font=\scriptsize] at (4.26,2.2) {$(2,4)$};
\node[font=\scriptsize] at (5.050000000000001,2.2) {$(3,7)$};
\node[anchor=east,font=\scriptsize] at (.61,1.66) {$(1,1)$};
\fill[black!16] (0.780,1.425) rectangle (1.420,1.895);
\node at (1.100,1.660) {$-$};
\fill[black!16] (1.570,1.425) rectangle (2.210,1.895);
\node at (1.890,1.660) {$-$};
\fill[black!16] (2.360,1.425) rectangle (3.000,1.895);
\node at (2.680,1.660) {$-$};
\node at (3.470,1.660) {$+$};
\node at (4.260,1.660) {$+$};
\node at (5.050,1.660) {$+$};
\node[anchor=east,font=\scriptsize] at (.61,1.0899999999999999) {$(1,2)$};
\node at (1.100,1.090) {$+$};
\node at (1.890,1.090) {$+$};
\node at (2.680,1.090) {$+$};
\fill[black!16] (3.150,0.855) rectangle (3.790,1.325);
\node at (3.470,1.090) {$-$};
\fill[black!16] (3.940,0.855) rectangle (4.580,1.325);
\node at (4.260,1.090) {$-$};
\node at (5.050,1.090) {$+$};
\node[anchor=east,font=\scriptsize] at (.61,0.52) {$(2,1)$};
\node at (1.100,0.520) {$+$};
\node at (1.890,0.520) {$+$};
\node at (2.680,0.520) {$+$};
\fill[black!16] (3.150,0.285) rectangle (3.790,0.755);
\node at (3.470,0.520) {$-$};
\node at (4.260,0.520) {$+$};
\fill[black!16] (4.730,0.285) rectangle (5.370,0.755);
\node at (5.050,0.520) {$-$};
\draw (.68,.215) rectangle (5.48,1.96);
\end{scope}
\begin{scope}[xshift=7.1cm]
\node at (2.75,2.65) {one choice of flips $A$};
\node[font=\scriptsize] at (1.1,2.2) {$(1,2)$};
\node[font=\scriptsize] at (1.8900000000000001,2.2) {$(2,3)$};
\node[font=\scriptsize] at (2.68,2.2) {$(3,4)$};
\node[font=\scriptsize] at (3.47,2.2) {$(1,3)$};
\node[font=\scriptsize] at (4.26,2.2) {$(2,4)$};
\node[font=\scriptsize] at (5.050000000000001,2.2) {$(3,7)$};
\node[anchor=east,font=\scriptsize] at (.61,1.66) {$(1,1)$};
\fill[black!16] (0.780,1.425) rectangle (1.420,1.895);
\node at (1.100,1.660) {$-$};
\node at (1.890,1.660) {$+$};
\draw (1.890,1.660) circle (.20);
\fill[black!16] (2.360,1.425) rectangle (3.000,1.895);
\node at (2.680,1.660) {$-$};
\node at (3.470,1.660) {$+$};
\node at (4.260,1.660) {$+$};
\node at (5.050,1.660) {$+$};
\node[anchor=east,font=\scriptsize] at (.61,1.0899999999999999) {$(1,2)$};
\node at (1.100,1.090) {$+$};
\node at (1.890,1.090) {$+$};
\node at (2.680,1.090) {$+$};
\node at (3.470,1.090) {$+$};
\draw (3.470,1.090) circle (.20);
\fill[black!16] (3.940,0.855) rectangle (4.580,1.325);
\node at (4.260,1.090) {$-$};
\node at (5.050,1.090) {$+$};
\node[anchor=east,font=\scriptsize] at (.61,0.52) {$(2,1)$};
\node at (1.100,0.520) {$+$};
\node at (1.890,0.520) {$+$};
\node at (2.680,0.520) {$+$};
\fill[black!16] (3.150,0.285) rectangle (3.790,0.755);
\node at (3.470,0.520) {$-$};
\node at (4.260,0.520) {$+$};
\node at (5.050,0.520) {$+$};
\draw (5.050,0.520) circle (.20);
\draw (.68,.215) rectangle (5.48,1.96);
\end{scope}
\draw[->,thick] (5.75,1.1)--(6.35,1.1);
\end{tikzpicture}
\caption{Three lines and six points from the scale-$3$ table. Rows index lines and columns index points. Shading marks negative entries; circles mark incidences flipped to positive. At the shared point $(1,3)$, the signs for lines $(1,2)$ and $(2,1)$ can differ. No two distinct lines share two incidence columns, so a negative rectangle must have a singleton side.}
\label{fig:incidence}
\end{figure}


\begin{proposition}[Geometry and an upper representation]\label{prop:geometry}
Every incidence flip satisfies
\begin{equation}
 \VC(\Hc_A)\le2,\qquad \rect(\Hc_A)\ge2^{-15},\qquad \dc(\Hc_A)\le N+3.
 \label{eq:geometry}
\end{equation}
The template has sign-rank at most six and has exactly
\begin{equation}
 I=2N^4-\bigl(N(N+1)/2\bigr)^2\ge N^4
 \label{eq:incidencecount}
\end{equation}
incidences.
\end{proposition}

\paragraph{Why the geometry survives.}
The score $(ax+b-y)^2-1/2$ gives a six-dimensional strict representation of $F$. The homogeneous-rectangle theorem of \citet{APP05}, in the formulation of \citet[Theorem~1.9]{HHPTZ22}, gives a large monochromatic rectangle under every product distribution. Positive rectangles survive every flip. A nonempty negative rectangle has a singleton side, because distinct lines meet at most once; splitting its other side by the new signs retains at least half its mass. This proves the rectangle bound. The same one-intersection property forbids shattering three points.

A representation needs to store only the labels along each line. Let $q_\ell(t)$ be its label at $(t,at+b)$, taking $+1$ when that point is outside the domain. Then
\begin{equation}
 q_\ell(x)+2(y-ax-b)^2
 \label{eq:representation-score}
\end{equation}
has the desired sign. It is linear in the $N+3$ features $\one_{\{x=t\}}$, $t\in[N]$, and $y,xy,y^2$. On the running line $y=x+1$, its score is exactly the chosen incidence sign. At every other grid point the squared residual is at least one, so the score is positive. Appendix~\ref{app:geometry} gives the coefficients and the incidence count.

\subsection{Counting on the incidence support}
The lower bound must count the part of the table that actually varies. Off-incidence entries contain no random information.

\begin{lemma}[Restricted sign patterns]\label{lem:mask}
For an integer $d\ge1$, on a specified set $J$ of $s$ entries of a $k\times q$ matrix, the number of strict sign patterns produced by real matrices of rank at most $d$ is at most
\begin{equation}
 \left(\frac{8es}{v}\right)^v,\qquad v=(k+q)d,\quad 1\le v\le s.
 \label{eq:maskcount}
\end{equation}
\end{lemma}

\begin{proof}
Factor the matrix into $k$ row vectors and $q$ column vectors in $\R^d$. Its entries in $J$ are $s$ degree-two polynomials in $v$ variables. Apply the strict-sign form of \citet{Warren68}; this factorization argument is standard in sign-rank counting \citep{AMY16}.
\end{proof}

For $0\le\eta<1/2$, let $\hbin(\eta)=-\eta\log_2\eta-(1-\eta)\log_2(1-\eta)$, with $0\log_2 0=0$, and set
\begin{equation}
 c_\eta=1-\hbin(\eta),\qquad
 s_\eta=\left\lfloor\frac{c_\eta^2 I}{32(K+Q)}\right\rfloor.
 \label{eq:seta}
\end{equation}
Lemma~\ref{lem:mask} counts at most $2^{c_\eta I/2}$ patterns in rank $s_\eta$; a Hamming ball of radius $\eta I$ contains at most $2^{\hbin(\eta)I}$ patterns. Their union covers at most a $2^{-c_\eta I/2}$ fraction of the random signs.

\begin{lemma}[Probabilistic reductions]\label{lem:prob-reduction}
If $\dc_\eta(\Hc_A)=d$, a strict matrix of rank at most $d+1$ makes at most $\eta I$ errors on the incidences. For every finite class,
\begin{equation}
 \dc_{1/4}(\Hc)\le\edc(\Hc)+1.
 \label{eq:exact-to-average}
\end{equation}
\end{lemma}

\paragraph{Why randomization does not evade the count.}
For each nonempty line support $J_\ell$, use the uniform marginal on $J_\ell$ and weight the row by $|J_\ell|/I$. The expected total incidence error is at most $\eta I$, so one embedding attains that bound. Make its predictions strict with the constant feature. For \eqref{eq:exact-to-average}, the constant feature provides a predictor of error at most $1/2$ whenever exact representation fails; failure has probability at most $1/2$. Appendix~\ref{app:counting} gives the full reductions.

\begin{theorem}[Dimension probabilities]\label{thm:dimensions}
A random incidence flip satisfies
\begin{align}
 \Pr[\dc(\Hc_A)\le s_0]&\le2^{-I/2},\label{eq:ordinaryprob}\\
 \Pr[\dc_\eta(\Hc_A)<s_\eta]&\le2^{-c_\eta I/2},\label{eq:averageprob}\\
 \Pr[\edc(\Hc_A)<s_{1/4}-1]&\le2^{-c_{1/4}I/2}.
 \label{eq:exactprob-bound}
\end{align}
\end{theorem}

Since $I\ge N^4$ and $K+Q=4N^3$, the thresholds are linear in $N$ for each fixed error. Proposition~\ref{prop:geometry} supplies the matching upper bound. Thus approximation and randomization leave the order of dimension unchanged for this family.

\section{Allowing a prior and discarding targets}\label{sec:prior}
A stronger relaxation chooses the feature law for a known target prior and requires it to work only for most targets. This is the average probabilistic dimension of \citet[Definition~3.2]{KM25}. We denote it by $\adc_{\eta,\delta}(\pi)$ and use Boolean threshold loss: it is the least $d$ admitting one law on maps $\phi:X\to\R^d$ with
\begin{equation}
 \Pr_{h\sim\pi}\left[\ \forall D\in\Delta(X),\quad
       \E_\phi\mathcal E_\phi(D,h)\le\eta\ \right]\ge1-\delta.
 \label{eq:adcdef}
\end{equation}
The law may depend on $\pi$. It may not depend on the realized target or on $D$. The good set of targets must work for every marginal.

\begin{theorem}[Dimension after discarding targets]\label{thm:prior}
Let $\Lc_0=[N]\times[N^2]$ be the full-length lines, let $L=|\Lc_0|=N^3$, and let $\pi_N$ be the prior on $\Hc_A$ induced by a uniform line in $\Lc_0$. Fix $0\le\eta<1/2$ and $0\le\delta<1$. Put
\begin{equation}
 g=\ceil{(1-\delta)L},\qquad
 t_{\eta,\delta}=\left\lfloor\frac{c_\eta^2 gN}{32(g+Q)}\right\rfloor.
 \label{eq:prior-threshold}
\end{equation}
Then
\begin{equation}
 \Pr_A[\adc_{\eta,\delta}(\pi_N)<t_{\eta,\delta}]
 \le \binom{L}{g}\,2^{-c_\eta gN/2}.
 \label{eq:prior-prob}
\end{equation}
Consequently $\adc_{\eta,\delta}(\pi_N)=\Theta_{\eta,\delta}(N)$ with probability $1-\exp(-\Omega_{\eta,\delta}(N^4))$.
\end{theorem}

\begin{proofsketch}
Fix $g$ rows. Their $gN$ incidence signs are independent. An embedding that approximates all these rows under every marginal gives, by testing on each row's line, an approximation of rank at most $d+1$ to these signs. Apply Lemma~\ref{lem:mask} with $g$ rows and $Q$ columns, then count nearby patterns. Finally take a union bound over the $\binom{L}{g}$ possible good row sets. Their logarithmic number is only $O(N^3)$; the random signs supply $\Theta(N^4)$ bits. Appendix~\ref{app:prior} proves the statement, including priors with repeated functions.
\end{proofsketch}

Even retaining only $1\%$ of the targets requires dimension linear in $N$. The constant depends on the retained fraction and the approximation error; the order does not.

The next proposition explains how the prior-average notion relates to a worst-target guarantee. It replaces the family of laws chosen separately for all priors by one common law, at a controlled increase in error.

\begin{proposition}[From all priors to all targets]\label{prop:prior-minimax}
For every finite class, $0\le\eta<1/2$, and $0\le\delta<1$, put $\theta=\eta+(1/2-\eta)\delta$. Then
\begin{equation}
 \dc_\theta(\Hc)-1
 \le \sup_{\pi\in\Delta(\Hc)}\adc_{\eta,\delta}(\pi)
 \le \dc_\eta(\Hc).
 \label{eq:prior-minimax}
\end{equation}
\end{proposition}

For the left inequality, append a constant feature. Good targets have expected error at most $\eta$; the remaining targets can use a constant predictor of error at most $1/2$. Hence every mixture of target--marginal pairs admits a feature law of average loss at most $\theta$. Minimax supplies one law for all pairs. The argument reduces feature maps to the finitely many Boolean prediction families they induce, rather than assuming compactness of real-valued maps. Appendix~\ref{app:prior} proves this reduction.

\subsection{Classical SQ dimension and a claimed representation bound}\label{sec:km}
The same construction separates prior-average representation from classical SQ dimension. Define $\sqdim(\Hc,D)$ as the largest number $d$ of distinct hypotheses satisfying $\E_D[h_i h_j]\le1/d$ for $i\ne j$, and put $\sqdim(\Hc)=\max_D\sqdim(\Hc,D)$. This is the one-sided correlation convention of \citet[Definition~3.5]{KM25}. Classical SQ dimension does not characterize distribution-independent strong learning \citep{Feldman17}.

\begin{proposition}[Rectangles bound classical SQ dimension]\label{prop:classical-sq}
For every finite nonempty class,
\begin{equation}
 \sqdim(\Hc)\le 2/\rect(\Hc).
 \label{eq:classical-sq}
\end{equation}
In particular, every incidence flip has $\sqdim(\Hc_A)\le2^{16}$.
\end{proposition}

On a monochromatic rectangle, the sum of $t$ selected hypotheses has squared magnitude $t^2$. The correlation bounds limit its expected square to $t+t(t-1)/d$, giving \eqref{eq:classical-sq}; see Appendix~\ref{app:prior}. Thus classical SQ dimension stays bounded while Theorem~\ref{thm:prior}'s prior-average dimension tends to infinity. No function of classical SQ dimension alone can bound the latter at fixed $\eta<1/2$ and $\delta<1$.

\citet[Theorem~4.1]{KM25} claim $\adc_{\eta,\delta}(\pi)=O(\sqdim(\Hc)^{24.01})$ for small fixed errors. \citet{FKS26} already report that the ``claimed upper bounds do not hold due to a flaw in the proof,'' citing a personal communication from Karchmer and Malach. Theorem~\ref{thm:prior} and Proposition~\ref{prop:classical-sq} give an explicit counterexample to Theorem~4.1 under its stated fixed-feature-law quantifiers, consistent with that report. The lower bound uses Boolean threshold loss and therefore also applies to their literal unthresholded $0$--$1$ loss, which is no smaller.

The proof of their mini-batch SGD theorem, Theorem~3.3, invokes Theorem~4.1. Our counterexample rules out that intermediate implication; it does not disprove the mini-batch SGD conclusion itself. Open Question~1 of \citet{FKS26} is separate. That question specifies an architecture and training procedure rather than formally defining a class of ``benign'' networks.

\section{Query optimality and general dimension bounds}\label{sec:optimality}
The learning upper bound holds for every flip, including the all-positive table, so a query lower bound cannot hold for every member. For random incidence signs, we now show that logarithmically many queries are necessary at fixed tolerance. The learner may know the entire realized table.

\begin{theorem}[Query lower bounds]\label{thm:querylower}
Let $\chi=1-\hbin(3/8)$. Suppose the integer $N\ge2$ satisfies $\chi N\ge6\log_2N$ (equivalently, $N\ge1373$). With probability at least $1-2^{-2N^3}-N^3e^{-N/32}$ over $A$, every randomized learner satisfying \eqref{eq:learnmodel}, for any $0<\eps<1/4$ and $0<\tau<1$, obeys
\begin{equation}
 m\ge\frac{\log_2 N+\log_2\!\bigl(\chi(1-8\eps/3)/8\bigr)}{\log_2\ceil{1/\tau}}.
 \label{eq:lowerimproper}
\end{equation}
If its output must belong to $\Hc_A$, then
\begin{equation}
 m\ge\frac{3\log_2 N+\log_2(1-3\eps)}{\log_2\ceil{1/\tau}}.
 \label{eq:lowerproper}
\end{equation}
The bounds already hold on targets from $\Lc_0$, with the uniform marginal on each target line.
\end{theorem}

\paragraph{Why outputs, not targets, must be counted.}
A fixed predictor fits a random row on its $N$ incidences with probability exponentially small in $N$. Counting all predictors shows that none fits more than $O(N^2)$ of the $N^3$ full-line instances. Quantizing the replies gives at most $\ceil{1/\tau}^m$ outputs for each fixed learner seed. Averaging over seeds proves \eqref{eq:lowerimproper}. A proper output fits at most one line instance, giving \eqref{eq:lowerproper}. An improper output can fit every line of one slope simultaneously, so the stronger proper bound cannot be obtained by counting target labels alone. Appendix~\ref{app:query} supplies both arguments.

The response tree also gives a converse bound on representation dimension. Its leaf predictors can serve as features, regardless of how the learner computes them. The next proposition extends this observation to all three dimension notions.

\begin{proposition}[Transcript representations]\label{prop:transcript}
Suppose a randomized depth-$m$ learner achieves expected error $\eps<1/2$ at tolerance $0<\tau<1$. Put $B_\tau=\ceil{1/\tau}$ and $\gamma=1-2\eps$. Then
\begin{align}
 \dc_\eps(\Hc)&\le B_\tau^m,\label{eq:transcript}\\
 \edc(\Hc)&\le B_\tau^m\ceil{2\gamma^{-2}\ln(2Q)},\label{eq:transcript-exactprob}\\
 \dc(\Hc)&\le B_\tau^m\ceil{2\gamma^{-2}\ln(2KQ)}.\label{eq:transcript-exact}
\end{align}
For a deterministic learner, $\dc(\Hc)\le B_\tau^m$.
\end{proposition}

For each learner seed, use all $B_\tau^m$ leaf outputs as coordinates. The rounded oracle selects one of them, proving \eqref{eq:transcript}. For exact representation, minimax converts the guarantee under every marginal into row-specific convex combinations with positive expected signed score at every point. Concatenating independent leaf maps and applying concentration proves \eqref{eq:transcript-exactprob}--\eqref{eq:transcript-exact}. Appendix~\ref{app:query} gives the full argument.

At fixed tolerance, the incidence separation attains the exponential order in $m$ allowed by these bounds. We do not claim that the logarithmic overhead for exact representations is necessary.

\begin{figure}[t]
\centering
\begin{tikzpicture}[x=1cm,y=1cm,>=Stealth,font=\small]
 \draw[->] (0,0)--(6.4,0) node[right,font=\footnotesize] {$n$};
 \draw[->] (0,0)--(0,3.35) node[above] {$\log_2 d$};
 \draw[dashed,thick] (.3,.17)--(5.8,3.0);
 \draw[thick] (.3,.1)--(5.8,2.32);
 \draw[thick] (.3,.18) .. controls (1.2,.51) and (2.9,.69) .. (5.8,.82);
 \node[anchor=west,align=left,font=\footnotesize] at (6.05,2.99)
   {transcript upper bound\\$m\log_2 B_\tau+O(\log n)=O_\eps(n)$};
 \node[anchor=west,align=left,font=\footnotesize] at (6.05,2.02)
   {incidence separation: $n/3+O(1)$\\$m=\Theta_\eps(n)$, $\tau=\Omega_\eps(1)$};
 \node[anchor=west,align=left,font=\footnotesize] at (6.05,.66)
   {Patel's lower bound: $\Omega(n^{1/3})$\\$m,\tau^{-1}=2^{\widetilde O(n^{2/9})}$};
\end{tikzpicture}
\caption{Growth orders of log dimension at fixed learning and representation errors; constants are suppressed. The incidence bound holds for all three dimension notions and meets the exponential order of the transcript bound at its own query resources. The dashed line uses $K=Q=2^n$, $m=\Theta_\eps(n)$, and fixed tolerance; it is not the upper bound at Patel's resources. Patel's curve denotes an ordinary-dimension lower bound, not an upper bound on his class.}
\label{fig:parameters}
\end{figure}


\section{Computational implementations}\label{sec:computation}
The statistical separation charges queries, not the size of the class description. Its computational content depends on whether the learner receives the full table or a short entry evaluator.

\begin{theorem}[Table and succinct implementations]\label{thm:table}
For every $0<\eps<1/4$, each incidence flip has the deterministic proper table-polynomial learner in \eqref{eq:mainlearn}. If its entries can instead be evaluated in time $E(n)$, it has a deterministic improper learner using $O(n^2)$ queries of tolerance $\eps/8192$, with computation polynomial in $(n,1/\eps,E(n))$ and an efficiently evaluable output.
\end{theorem}

The table algorithm derives a constant-size catalog of geometric rectangles from the six-dimensional template, then chooses useful rectangles deterministically. At the first negative labelling step, the one-intersection property either identifies a good row or locates a negative point. The remaining row supports away from that point are disjoint, which permits proper completion. Appendix~\ref{app:proper} gives a finite-bit implementation.

The succinct algorithm uses the order along a line. At a trial slope $a$, the residual $y-ax$ increases with $x$ if $a$ is too small and decreases if it is too large. Noisy medians and one intersection-mass query distinguish these cases. A heavy atom identifies either the whole line or one negative point. This yields $O(\log N)$ stages with $O(\log N)$ queries each. Appendix~\ref{app:median} proves the result and a proper variant under a condition satisfied by almost every random table.

A succinct separation follows under a cryptographic assumption. Under the direct subexponential pseudorandom-function assumption stated in Appendix~\ref{app:prf}, for every prescribed integer $c\ge1$ there are polynomial-length keys such that almost every key defines a class with all three dimensions at least $n^c$ and a key-known randomized proper polynomial-time learner using $O_\eps(n^2)$ queries at fixed tolerance. Every key has the deterministic improper learner. The key-length exponent depends on $c$; this does not give an unconditional succinct construction, or an exponential gap for one polynomial-time family.

\section{Discussion}\label{sec:open}
The separation is sharp in dimension for the incidence construction and in query order at fixed tolerance. Its exponential gap still relies on an unstructured table: the counting argument does not choose a high-dimension flip with an unconditional polynomial-time evaluator. The conditional succinct families give each prescribed polynomial dimension separately, not this exponential gap. The full dependence on a varying tolerance remains undetermined. None of these SQ results resolves the separate gradient-descent question.

\ifreview\else
\acks{GPT-6 Astra and GPT-6 Astra Pro (OpenAI) assisted with mathematical development, source comparison, code, and writing. Claude (Anthropic) assisted with the Lean formalization. The author is responsible for the results and presentation.}
\fi
\paragraph{Reproducibility.}
The source includes complete proofs, build instructions, and finite arithmetic checks. The Lean development checks 213 supporting theorems for incidence geometry, the classical-SQ bound, noisy medians, conditional counting and query bounds, and the rectangle learner's proof steps. Its coverage table records the external hypotheses and the parts not yet formalized; no whole-paper formal verification is claimed.
\label{main-end}
